\renewcommand{\bibname}{Quellen und Bandzuordnung}
\begin{thebibliography}{9}
\addcontentsline{toc}{chapter}{Quellen und Bandzuordnung}
\footnotesize
\setlength{\itemsep}{4pt}
\setlength{\parskip}{0pt}
\interlinepenalty=10000

\bibitem{project}
Martin Kunze: \emph{Die Grundlagen der Mathematik}.
Band~8: Begriff der Bijektion; Einschränkungen und Korestriktionen;
Umkehrfunktion; Kompositionen und Inversen; induzierte
Potenzmengenabbildung; Verkleben disjunkter Zweige;
Erweiterung um ein neues Paar; Bijektionen mit vorgeschriebenem Wert;
leere Bereiche. Band~3 und Band~5 liefern Mengen- und Funktionsnotation.
Die Darstellung folgt den beweisfreien Leitfäden zu Band~3,
Band~4, Band~5, Band~28 und Band~37.

\bibitem{hammack}
Richard Hammack: \emph{Book of Proof}.
3.~Auflage, Fassung~3.4.
Abschnitt~12.2, Druckseite~229: Begriff der Bijektion.
Abschnitt~12.5, Satz~12.3, Druckseite~240:
Bijektivität genau dann, wenn der umgedrehte Graph eine
Funktion von der Zielmenge in den Definitionsbereich ist.
Abschnitt~12.5, Druckseiten~240--241: beidseitige Umkehrgleichungen.
\url{https://richardhammack.github.io/BookOfProof/Main.pdf}.

\end{thebibliography}

\noindent\textbf{Zuordnung und Auswahl.}
Die Charakterisierungen, Umkehrgesetze und Ergebnisse über
Teilmengen, Einschränkungen und Erweiterungen fassen die genannten
Themen aus Band~8 zusammen. Die Kapitel über Produkte,
Umindizierung und Funktionenmengen ergänzen diese durch
mengentheoretische Standardkonstruktionen mit ausdrücklich
angegebenen Umkehrungen. Der Zweipunkttest ist eine zusätzliche
Charakterisierung auf einer Funktionsmenge.

\medskip
\noindent\textbf{Grenzen der Zusammenstellung.}
Die Formulierungen und ihre Zusammenfassung sind eigenständig;
nicht jede Aussage steht in den Quellen in genau dieser Gestalt.
Es werden weder historische Erstentdeckungen noch eine vollständige
Erfassung aller bekannten Resultate beansprucht.
Reine Injektivitäts- und Surjektivitätssätze stehen in Band~6 und
Band~7; allgemeine Funktionssätze in Band~5. Auswahlprinzipien,
Äquivalenzrelationen, Ordnungen und zusätzliche algebraische
Strukturen sind keine Gegenstände dieses Leitfadens.
